\chapter{ЗАКЛЮЧЕНИЕ}

В ходе работы разработан многослойный персептрон с настраиваемой архитектурой и обучен
решению задачи бинарной классификации диагноза. Реализованы
прямое распространение сигнала, бинарная кросс-энтропия и две другие функции потерь,
обратное распространение ошибки и градиентный спуск с возможностью накопления импульса.

Обратный проход построен на обратном автоматическом дифференцировании: граф вычислений
формируется как побочный продукт прямого прохода, а производные по параметрам выводятся
программно из правил для элементарных операций. 

Архитектура выбрана экспериментально сравнением шести конфигураций, усреднённых по пяти
запускам. Установлено, что в рассматриваемой задаче ёмкость сети не является
ограничивающим фактором: изменение числа параметров в двенадцать раз не приводит к
значимому изменению результата, а добавление третьего скрытого слоя его ухудшает.
Определяющим параметром оказалась функция активации скрытых слоёв: замена сигмоиды на
гиперболический тангенс снижает функцию потерь на проверочной выборке втрое при
неизменном числе параметров, что объясняется четырёхкратной разницей в максимуме
производной и симметричностью тангенса относительно нуля. Функция \texttt{ReLU} при этом
даёт наилучшее приближение обучающей выборки, но заметно переобучается и показывает
втрое больший разброс по сиду генератора случайных чисел.

В качестве итоговой принята конфигурация из двух скрытых слоёв по 24 нейрона с
гиперболическим тангенсом, содержащая 1394 параметра. На проверочной выборке она даёт
бинарную кросс-энтропию 0,0111; усреднение
по пяти сидам даёт 0,0265 и 0,9858 соответственно, при постоянном ответе 0,6283.
Признаков переобучения не наблюдается: разрыв между функциями потерь на обучающей и
проверочной выборках сокращается на протяжении всего обучения.

